\section{Stability Across Eras}
\label{sec:era_stability}

Every attribution result in this paper is a full-sample aggregate over twenty
years. That is the convention in this literature and it hides something
important on this panel.

\subsection{The exogenous block helps in seven eras of eight}

\begin{table}[H]
\centering
\caption{Exogenous designs against the backbone, by era, under the rolling
protocol. Negative means the design beats the backbone. The stripped-down
amplitude design is included because its behaviour in the failing era is the
point---and because that behaviour turned out to be an artefact of the data
defect of Section~\ref{sec:stale_data} rather than a property of the design.}
\label{tab:eras}
\begin{tabular}{llrrr}
\toprule
Block & Period & backbone & ridge on 492 & stripped-down \\
\midrule
1 & 2007-02--2009-04 & 0.15054 & $-0.00400$ & $-0.00271$ \\
2 & 2009-04--2011-06 & 0.10915 & $-0.00222$ & $+0.00526$ \\
3 & 2011-06--2013-08 & 0.11269 & $-0.00231$ & $-0.00105$ \\
4 & 2013-08--2015-09 & 0.13143 & $-0.00265$ & $-0.00258$ \\
5 & 2015-09--2017-11 & 0.12048 & $-0.00469$ & $-0.00323$ \\
6 & 2017-11--2020-01 & 0.15594 & $-0.00591$ & $-0.00510$ \\
7 & 2020-01--2022-03 & 0.17033 & $-0.00771$ & $-0.00977$ \\
8 & 2022-03--2024-04 & 0.13511 & $\mathbf{+0.03321}$ & $+0.00347$ \\
\bottomrule
\end{tabular}
\end{table}

The exogenous block beats the backbone in seven consecutive eras---including
through 2008 and 2020---by $0.0022$ to $0.0077$, and then fails once by
$+0.03321$, an order of magnitude larger than any era's gain. That single block
erases seven eras of advantage: over the whole panel the raw exogenous design
sits at $0.13617$ against the backbone's $0.13571$, a difference of $+0.00046$
that is indistinguishable from zero.

A twenty-year aggregate therefore reports ``exogenous information helps'' or
``exogenous information does not help'' depending on where the window ends,
and both statements are true of their own windows. We regard the aggregate
alone as insufficient reporting, in this paper and in the literature it
follows.

\subsection{The failure is 27 bars, and it is our own}

Concentration within the failing era, over its $27{,}376$ bars:

\begin{table}[H]
\centering
\caption{Share of the era's total loss differential carried by its worst bars.}
\label{tab:concentration}
\begin{tabular}{lrr}
\toprule
Worst & Bars & Share of damage \\
\midrule
$0.1\%$ &    27 & $47.2\%$ \\
$1.0\%$ &   273 & $90.5\%$ \\
$5.0\%$ & 1{,}368 & $122.4\%$ \\
$10.0\%$ & 2{,}737 & $138.0\%$ \\
\bottomrule
\end{tabular}
\end{table}

Shares above $100\%$ mean the remaining bars have a \emph{negative}
differential: on the typical bar of the failing era the exogenous block still
helps, and the median per-bar differential is $+0.00069$ against a mean of
$+0.03318$. Twenty-seven bars out of twenty-seven thousand carry half of it.

All ten worst bars fall on 13--15 October 2023, and one channel,
\texttt{voldemand\_all\_open\_only}, carries $97.2\%$ of the negative swing.
Zero of forty-one channels exceeded three times their prior-month range, so
this is not a market event. It is the imputation defect of
Section~\ref{sec:stale_data}: a series absent from the raw data, forward-
filled, scaled by a collapsed interquartile range into coordinates in the
hundreds, and fitted a coefficient.

Two things follow. The reported \texttt{vol\_demand} bucket effect
(Section~\ref{sec:multiplicity_results}) rests partly on that artifact and is
suspended pending re-estimation on observed bars. And the practical reading of
the exogenous block is not that it stopped working---by 2024 the stripped-down
design is helping again---but that carrying it exposes a forecaster to rare
collapses that require detection rather than tolerance.
